%!TEX root = ./../../thesis-trier.tex

\IN{bottom-up method} 
\IN{extension method}
\IN{augmentative subset}
Let $\semiautomat{A}{Q}{\Sigma}{\delta}$. A subset $W \subseteq \Sigma^*$
is called \emph{augmentative (for $\mathcal A)$} if for every non-empty proper subset $S \subseteq Q$
there exists $u \in W$ such that $|\applytrans{\delta^{-1}}{S}{u}| > |S|$. Let $W \subseteq \Sigma^*$ be augmentative.
Since the number of subsets of states is finite, 
every augmentative subset contains a finite augmentative subset.
Hence, we can assume~$W$ is finite.
Clearly, if~$\mathcal A$ admits an augmentative subset,
then it is synchronizing
and there exist $q \in Q$ and $a \in \Sigma$
so that $|\applytrans{\delta^{-1}}{\{q\}}{a}| \ge 2$. Then, by using $W$ to successively enlarge
this subset, we obtain the following bound (see also~\cite[Thm.~2]{Rystsov95a}):
\begin{align}
		\rt{\mathcal A} 
		& \le \max\{ |u| : u \in W \} \cdot (n - |\applytrans{\delta^{-1}}{\{q\}}{a}|) + 1 \nonumber \\
		& \le \max\{ |u| : u \in W \} \cdot (n - 2) + 1. \label{sync_max:eqn:rt_bound}
\end{align}
Obtaining bounds for the reset length by using an augmentative set with \cref{sync_max:eqn:rt_bound}
is called the \emph{bottom-up} (or \emph{extension}) \emph{method}~\cite{DBLP:journals/ijfcs/Berlinkov11}. 

The next statement from Rystsov~\cite{Rystsov95a,Rystsov2000}
characterizes automata admitting an augmentative subset
as the strongly connected synchronizing automata.

\begin{theorem}[Rystsov~\cite{Rystsov95a,Rystsov2000}]
	\label{sync_max:thm:augmentative}
	Let $\semiautomat{A}{Q}{\Sigma}{\delta}$.	
	Then $\mathcal A$ is strongly connected and synchronizing
	if and only if $\mathcal A$ admits an augmentative subset $W \subseteq \Sigma^*$.	
\end{theorem}
\begin{proof}
    First, assume $\mathcal A$ is strongly connected and synchronizing.
	Then for every $q \in Q$ there exists $u_q \in \Sigma^*$
	such that $\applytrans{\delta}{Q}{u_q} = \{q\}$.
	In this case, since for a proper and non-empty $S \subseteq Q$
	and $q \in S$ we have $|\applytrans{\delta^{-1}}{S}{u_q}| = |Q| > |S|$,
	the set $W = \{ u_q \mid q \in Q\}$	is augmentative.
	Conversely, assume $W \subseteq \Sigma^*$
	is augmentative. Then, every state $q \in Q$ is synchronizing
	as we can enlarge $\{q\}$ successively to obtain the whole state set.
	So, by Corollary~\ref{prelim:cor:SCC_sync},
	$\mathcal A$ is strongly connected.
\end{proof}

Recall that, for an OI-automaton $\mathcal A$,~$\Sigma_I$
denotes the set of idempotent one-point singular letters and~$\Sigma_G$ the set of permutational letters.
Rystsov~\cite[Lem.~2]{Rystsov2000} has shown the following for SI-automata.
The proof works for OI-automata without modification. We give it for the convenience of the reader.

\begin{proposition} \label{sync_max:prop:SI-augmentative-set-form}
	Let $\semiautomat{A}{Q}{\Sigma}{\delta}$
	be an OI-automaton and $S \subseteq Q$.
	If $u \in \Sigma^*$ is a shortest word such
	that $|S| < |\applytrans{\delta^{-1}}{S}{u}|$,
	then $u \in \Sigma_I\Sigma_G^*$.	
\end{proposition}
\begin{proof}
	We can write $u = vaw$
	with $v \in \Sigma^*$, $a \in \Sigma_I$ and $w \in \Sigma_G^*$.
	Assume $|v| \ge 1$.
	Then, by the minimal choice of $u$,
	we have $|\applytrans{\delta^{-1}}{S}{aw}| \le |S|$.
	Set $T = \applytrans{\delta^{-1}}{S}{w}$.
	As~$w$ permutes $Q$, we have $|T| = |S|$.
	By \cref{prelim:eqn:delta_inverse}
	we have $\delta^{-1}(S, aw) = \delta^{-1}(T, a)$. Thus we find
	$|\applytrans{\delta^{-1}}{T}{a}| \le |T|$.
	Hence $T$ cannot be properly contained in $\applytrans{\delta^{-1}}{T}{a}$.
	By~\crefitem{prelim:thm::one-attractor-idempotent}{prelim:prop:item:inverse_super_or_subset2},
	this gives $\applytrans{\delta^{-1}}{T}{a} \subseteq T$.	
	Then, with \cref{prelim:eqn:delta_inverse}, we find
	\[
	 \applytrans{\delta^{-1}}{S}{vaw}
	 = \applytrans{\delta^{-1}}{\applytrans{\delta^{-1}}{T}{a}}{v}
	 \subseteq \applytrans{\delta^{-1}}{T}{v} = \applytrans{\delta^{-1}}{S}{vw},
	\]
	which gives $|S| < |\applytrans{\delta^{-1}}{S}{vw}|$, 
	contradicting the minimal choice of $u$. So $v = \varepsilon$ and this gives $u \in \Sigma_I\Sigma_G^*$.
\end{proof}

The next corollaries we give by Rystsov\cite{Rystsov2000,Rystsov2022} 
for SI-automata, where \cref{sync_max:cor:aug:bounds:semiregular} was only stated
for the case that the group of units is regular.

\begin{corollary}
	\label{sync_max:cor:augmentative_subset_IG}
	Let $\semiautomat{A}{Q}{\Sigma}{\delta}$ be an OI-automaton.
	If there exists an augmentative subset,
	then there exists an augmentative subset in~$\Sigma_I \Sigma_G^*$. 
\end{corollary}
\begin{proof}
	For each non-empty and proper subset $S \subseteq Q$
	let $u_S \in \Sigma^*$ be a shortest
	words such that $|\delta^{-1}(S, u_S)| > |S|$.
	The set of all these words $u_S$ for some non-empty and proper $S \subseteq Q$
	is augmentative and, by \cref{sync_max:prop:SI-augmentative-set-form},
	a subset of $\Sigma_I \Sigma_G^*$.	
\end{proof}

The next bound also follows from \cref{prelim:prop:SI-without-perm:reach-bound} for SI-automata
and is tight even for SI-automata.

\begin{corollary} \label{syncmax:graph:cor:OI-sync}
	Let $\semiautomat{A}{Q}{\Sigma}{\delta}$ be an $n$-state synchronizing OI-automaton without permutational letters.
	Then a shortest synchronizing word has length at most $n - 1$.
\end{corollary}
\begin{proof}
	Since~$\mathcal A$ is synchronizing, there exist $q \in Q$ and $u \in \Sigma^*$
	such that $\delta^{-1}(q, u) = Q$.
	Set $S_1 = \{q\}$ and, inductively, construct a sequence $S_1, \ldots, S_m$ with $m < |Q|$
	such that $S_1 \subseteq \cdots \subseteq S_m$ with $S_m = Q$
	and where, for each  $i \in \{2,\ldots,m\}$, 
	\begin{enumerate}
		\item $S_{i} = \delta^{-1}(S_{i-1}, a_i)$ for some $a_i \in \Sigma_I$ and
			
		\item $|S_i| > |S_{i-1}|$	
	\end{enumerate}
	in the following way: Let $i \in \{2,\ldots,m\}$. Let $i \in \{2,\ldots,m\}$. Since $q \in S_{i-1}$
	and $S_{i-1} \ne Q$ we have $|\delta^{-1}(S_{i-1}, u)| > |S_{i-1}|$.
	Thus there exists a shortest $v \in \Sigma^*$
	such that $|\delta^{-1}(S_{i-1}, v)| > |S_{i-1}|$.
	By \cref{sync_max:prop:SI-augmentative-set-form} 
	we have $v \in \Sigma_I \Sigma_G^*$.
	Since $\Sigma_G$ is empty, this yields $v = a_i$
	for some $a_i \in \Sigma_I$.
	Set $S_i = \delta^{-1}(S_{i-1}, a_i)$.
	By~\crefitem{prelim:thm::one-attractor-idempotent}{prelim:prop:item:inverse_super_or_subset2},
	we find $S_{i-1} \subseteq S_i$.
	By finiteness the sets $S_i$ cannot grow indefinitely. Thus, 
	this process must come to a halt after at most~$m$ steps with $m < |Q|$
	and $S_m = Q$.
	The word $a_m \cdots a_1$ synchronizes $\mathcal A$.
\end{proof}

\begin{corollary} \label{sync_max:cor:aug:bounds:semiregular} 
	Let $\semiautomat{A}{Q}{\Sigma}{\delta}$ be a strongly connected and synchronizing OI-automaton.
	If the group of units of $\mathcal A$ is a semiregular permutation group,
	then $\rt{\mathcal A} \le (n-1)^2$ and this bound is tight.
\end{corollary}
\begin{proof}
	By \cref{sync_max:prop:grp_of_units_TA} for an OI-automaton the group of 
	units is $G = \{ \fword{\delta}{u} \mid u \in \Sigma_G^* \}$.
	For $g \in G$ let $u \in \Sigma_G^*$ be a shortest permutational word
	such that $g = \fword{\delta}{u}$.	
	Assume $|u| \ge n$. 
	Then $\delta(q, u_1) = \delta(q, u_1 u_2) = \delta(\delta(q, u_1), u_2)$
	for some $u_1, u_2, u_3 \in \Sigma_G^*$
	with $u = u_1 u_2 u_3$, $|u_2| > 0$, and $|u_1| < n$. So $\fword{\delta}{u_2}$
	has the fixpoint $\delta(q, u_2)$,
	which gives $\fword{\delta}{u_1} = \operatorname{id}_{Q}$
	as $G$ is semiregular.
	Thus $g = \fword{\delta}{u_1 u_2}$, contradicting the minimal choice of~$u$.
	Thus $|u| < n$, which gives $G = \{ \fword{\delta}{u} \mid u \in \Sigma_G^{< n} \}$.
	This yields that every transformation induced by a word in $\Sigma_I \Sigma_G^*$
	is also induced by a word in $\Sigma_I \Sigma_G^{< n}$,
	and the claim follows from \cref{sync_max:cor:augmentative_subset_IG}
	and \cref{sync_max:eqn:rt_bound}. Tightness of this bound
	follows from the \v{C}ern\'y automata (\cref{intro:fig:cerny})
	whose group of units is regular.
\end{proof}